\providecommand{\Res}[2]{\qty[round-mode=figures,round-precision=3]{\fpeval{#1}}{#2}}
\providecommand{\Val}[1]{\num[round-mode=figures,round-precision=4]{\fpeval{#1}}}
\subsubsection*{Question a) Hauteur initiale}
\begin{enumerate}[label=\arabic*)]
    \item \textbf{Variables connues et recherchées :}
    \begin{itemize}[label= -]
        \item $L = \qty{\v{L}}{m}$, $\theta = \ang{\v{th}}$
        \item Variable recherchée : $h$ (m)
    \end{itemize}
    \item \textbf{Équations pertinentes :} la boule est à une distance verticale $L\cos\theta$ sous le point d'attache, au lieu de $L$ :
    \[ h = L - L\cos\theta \]
    \item \textbf{Isolation de la variable recherchée :} \[ h = L(1 - \cos\theta) \]
    \item \textbf{Résolution numérique :}
    \[ h = \num{\v{L}} \cdot (1 - \cos(\ang{\v{th}})) = \Res{\v{L}*(1-cosd(\v{th}))}{m} \]
    \item \textbf{Réponse :} la boule est soulevée de \textbf{\Res{\v{L}*(1-cosd(\v{th}))}{m}}.
\end{enumerate}

\subsubsection*{Question b) Vitesse au point le plus bas}
\begin{enumerate}[label=\arabic*)]
    \item \textbf{Variables connues et recherchées :}
    \begin{itemize}[label= -]
        \item $h = \Res{\v{L}*(1-cosd(\v{th}))}{m}$, $v_0 = \qty{0}{m/s}$
        \item Variable recherchée : $v$ (m/s)
    \end{itemize}
    \item \textbf{Équations pertinentes :} \[ m g h = \tfrac{1}{2} m v^2 \]
    \item \textbf{Isolation de la variable recherchée :} \[ v = \sqrt{2 g h} \]
    \item \textbf{Résolution numérique :}
    \[ v = \sqrt{2 \cdot \num{9.81} \cdot \Val{\v{L}*(1-cosd(\v{th}))}} = \Res{sqrt(2*9.81*\v{L}*(1-cosd(\v{th})))}{m/s} \]
    \item \textbf{Réponse :} au point le plus bas, la boule passe à \textbf{\Res{sqrt(2*9.81*\v{L}*(1-cosd(\v{th})))}{m/s}}.
\end{enumerate}

\subsubsection*{Question c) Énergie cinétique au point le plus bas}
\begin{enumerate}[label=\arabic*)]
    \item \textbf{Variables connues et recherchées :}
    \begin{itemize}[label= -]
        \item $m = \qty{\v{m}}{kg}$, $h = \Res{\v{L}*(1-cosd(\v{th}))}{m}$
        \item Variable recherchée : $E_c$ (J)
    \end{itemize}
    \item \textbf{Équations pertinentes :} toute l'énergie potentielle est devenue cinétique :
    \[ E_c = m g h \]
    \item \textbf{Isolation de la variable recherchée :} l'équation est déjà isolée.
    \item \textbf{Résolution numérique :}
    \[ E_c = \num{\v{m}} \cdot \num{9.81} \cdot \Val{\v{L}*(1-cosd(\v{th}))} = \Res{\v{m}*9.81*\v{L}*(1-cosd(\v{th}))}{J} \]
    \item \textbf{Réponse :} l'énergie cinétique est de \textbf{\Res{\v{m}*9.81*\v{L}*(1-cosd(\v{th}))}{J}}.
\end{enumerate}

\subsubsection*{Question d) Vitesse à \ang{\fpeval{\v{th}/2}}}
\begin{enumerate}[label=\arabic*)]
    \item \textbf{Variables connues et recherchées :}
    \begin{itemize}[label= -]
        \item $L = \qty{\v{L}}{m}$, $\theta_0 = \ang{\v{th}}$, $\theta_1 = \ang{\fpeval{\v{th}/2}}$
        \item Variable recherchée : $v_1$ (m/s)
    \end{itemize}
    \item \textbf{Équations pertinentes :} \[ m g L(1 - \cos\theta_0) = m g L(1 - \cos\theta_1) + \tfrac{1}{2} m v_1^2 \]
    \item \textbf{Isolation de la variable recherchée :} \[ v_1 = \sqrt{2 g L (\cos\theta_1 - \cos\theta_0)} \]
    \item \textbf{Résolution numérique :}
    \[ v_1 = \sqrt{2 \cdot \num{9.81} \cdot \num{\v{L}} \cdot (\cos(\ang{\fpeval{\v{th}/2}}) - \cos(\ang{\v{th}}))} = \Res{sqrt(2*9.81*\v{L}*(cosd(\v{th}/2)-cosd(\v{th})))}{m/s} \]
    \item \textbf{Réponse :} à cet angle, la boule passe à \textbf{\Res{sqrt(2*9.81*\v{L}*(cosd(\v{th}/2)-cosd(\v{th})))}{m/s}}.
\end{enumerate}
